\section{Heat cells and continuous color scales}
\label{sec:heat}
Heat encodes one scalar through fill color and retains its formatted value inside
the rectangle. Use \key{\gradheat[options]{value}} in text, ordinary tables or
existing TikZ pictures. Color is interpolated linearly over min/max.
\begin{demo}{Continuous colors and a neutral missing cell}
\gradheat[min=0,max=100]{0}\quad
\gradheat[min=0,max=100]{50}\quad
\gradheat[min=0,max=100]{100}\quad
\gradheat{NA}
\end{demo}
\subsection{Colors, text and missing observations}
\begin{sloppypar}
\key{heat low color} and \key{heat high color} specify the endpoint colors;
better=lower reverses their direction. \key{heat text color=auto} chooses black
or white by contrast, or accepts an explicit color. \key{heat missing color}
provides a neutral missing fill, distinct from zero. Precision, units, number
formatting, custom text, rounded corners, print grayscale and alternative text
work. Labels are inside; \key{label=none} hides a standalone cell's label. Bar targets,
errors, patterns and threshold bins are not drawn.
\end{sloppypar}
\subsection{Diverging colors and a center value}
\key{heat palette=diverging} selects orange/white/teal. Heat center defaults to
0 and heat center color to white; the center may instead be a target value.
It must lie within min/max, including their boundaries.
\begin{demo}{Zero and target centers in asymmetric ranges}
\gradheat[heat palette=diverging,
  min=-50,max=100,heat center=0]{-25}\quad
\gradheat[heat palette=diverging,
  min=-50,max=100,heat center=0]{0}\quad
\gradheat[heat palette=diverging,
  min=-50,max=100,heat center=0]{50}

\gradheat[heat palette=diverging,
  min=20,max=80,heat center=50]{35}\quad
\gradheat[heat palette=diverging,
  min=20,max=80,heat center=50]{50}\quad
\gradheat[heat palette=diverging,
  min=20,max=80,heat center=50]{65}
\end{demo}
Each side interpolates independently, so the center stays neutral in asymmetric
ranges. Better=lower swaps endpoints, retaining the center color. Print uses
gray endpoints and a white center. Place custom heat low/center/high color keys
after the palette; \key{heat palette=sequential} retains the earlier blue scale.
Automatic ranges do not expand to include a center outside the data range.
Keep \key{include zero=true} for a zero center, or set appropriate fixed limits.
Missing retains its separate fill; raw labels remain unchanged and alternative
text also describes the center.

\subsection{Declaring a heat column in gradtable}
Use \key{type=heat} with one scalar or NA per cell. Fill covers content width,
retaining column spacing. Best/second-best marking uses \key{mark=best|both}.
This example uses a fixed range; see Section~\ref{sec:column-scale} for automatic
scales and Section~\ref{sec:csv-input} for CSV data sources.
\begin{demo}{A fixed-scale heat column}
\begin{gradtable}[width=\linewidth,header align=c]
  \gradcolumn{name}{Method}
  \gradcolumn[type=heat,mark=both,
    options={min=0,max=100,precision=0}]{score}{Score}
  \gradrow{A,70}
  \gradrow{B,90}
  \gradrow{Pending,NA}
\end{gradtable}
\end{demo}
